\chapter{Literature Review}
\label{chap:related_work}

This chapter locates the thesis at the intersection of misinformation diffusion,
online-firestorm theory, geoengineering discourse, network science, and
large-language-model (LLM) agent simulation.  The distinction between these
literatures matters.  A model may reproduce the reach of a cascade without
representing how its claims change, while fluent agent output may look socially
plausible without constituting a valid model of people.  The review therefore
connects substantive theory to the thesis's observables and closes by stating
the specific gap addressed by the experiments.

\section{Misinformation as diffusion and semantic transformation}
\label{sec:lit-misinformation}

Misinformation is not a single technical object.  The term can refer to false
content, an actor's belief, an exposure, a sharing decision, or a population-level
diffusion pattern.  \textcite{lazer2018science} consequently call for research
that joins individual cognition, institutions, and platform structure rather
than treating falsity as a classification problem alone.  Large-scale
observational evidence confirms that diffusion is socially structured:
\textcite{vosoughi2018false} found that false news on Twitter travelled farther,
faster, deeper, and more broadly than true news in their data.  Novelty and
emotional response helped distinguish false from true cascades, while automated
accounts accelerated both.  These findings motivate models with content and
agent state, but they do not imply that every simulated cascade should reproduce
the same ordering.

Classical diffusion models provide useful null formalisms.  In the independent
cascade model, an activated node receives one probabilistic opportunity to
activate each neighbour \parencite{kempe2003influence}.  Related compartmental
and threshold models reduce the state of a node to susceptible, active, adopted,
or an opinion scalar.  Their abstraction is an advantage for identifying
conditions for reach or consensus, but an edge transmits a state rather than
rewriting a proposition.  This thesis instead studies a transmission in which
the text itself can lose, invert, or acquire claims.  The distinction is between
\emph{whether} something propagates and \emph{what it becomes while propagating}.

Empirical cascade size is also not synonymous with virality.
\textcite{goel2016virality} show that similarly popular diffusion events range
from broadcast-like stars to multigenerational chains and define structural
virality from pairwise distances in a diffusion tree.  A scale-free contagion
model did not recover the full empirical diversity they observed.  This is a
warning against explaining an outcome from a topology label alone.  In the
present study, event count, semantic degradation, and time to a severity
threshold are therefore conceptually separate outcomes.

\section{Online firestorms: theory and operationalisation}
\label{sec:lit-firestorms}

\textcite{pfeffer2014firestorms} define an online firestorm as a sudden discharge
of large quantities of negative word-of-mouth and complaint behaviour directed
at a person, organisation, or group.  Their contribution is a qualitative
synthesis of the mechanisms that make rapid, affective outrage possible.  The
paper's outlook identifies \emph{seven} interrelated factors: speed and volume
of communication, binary choices, network clusters, unrestrained information
flow, lack of diversity, cross-media dynamics, and network-triggered decision
processes.  It does not present an agent-based generative model, a semantic
misinformation scale, or a tipping-time estimator.

This point prevents a common attribution error.  The six simulation dimensions
used in this thesis---valence, surprise, identity alignment, network clustering,
information echo, and temporal acceleration---are a \emph{computational
remapping}, not Pfeffer et al.'s original list.  Valence represents the affective
character of firestorm messages.  Identity alignment and information echo
approximate lack of diversity and clustered reinforcement.  Network clustering
maps directly to local structure, while surprise and temporal acceleration
separate seeding shocks from the speed of subsequent exchange.  Cross-media
dynamics has no equivalent in the present single-platform society and is a
boundary condition.  Nor does the simulator reproduce the original paper's
binary platform affordances in a literal way.  The remapping is thus an
operational hypothesis to be assessed, not a claim that all seven mechanisms
have been implemented.

Later work has moved toward measuring observed firestorms.
\textcite{nuortimo2020scale} propose a three-level scale based on a media-analytics
case, and \textcite{nuortimo2025scale} develop dimensions of firestorm width,
intensity, and peak duration with AI-assisted validation.  These studies make
event severity more explicit, but remain different from a claim-level instrument
applied during a controlled generative cascade.  The literature leaves open how
a theory of rapid collective outrage should be connected to the semantic
mutation of scientific information.

\section{Geoengineering discourse as the substantive case}
\label{sec:lit-geoengineering}

Solar geoengineering is an unusually demanding case for studying that
connection.  Discussion combines scientific uncertainty, climate politics,
governance, distrust, and established conspiracy narratives.  It therefore
cannot be represented adequately as generic ``fake news.''  The principal
empirical grounding is \textcite{debnath2023conspiracy}, who analysed 814,924
English-language tweets containing \texttt{\#geoengineering} from 2009 to 2021.
Using natural-language processing, toxicity analysis, word embeddings, and
network analysis, they show that the false chemtrails narrative substantially
shaped geoengineering discussion.  Conspiracies spilled into debates in the
United Kingdom, United States, India, and Sweden and connected with broader
political themes.  Positive emotions tended to increase following governance
events, whereas negative and neutral emotions increased around project and
experiment announcements; toxicity widened anti-solar-radiation-management
spillovers.

Debnath et al. establish an observational association, not a causal diffusion
model or an LLM-agent society.  Their corpus is nevertheless valuable for
defining plausible discourse tendencies and seed events.  This thesis derives
reduced belief profiles for conspiracy, climate-action, and environmental-concern
positions from that account.  These profiles are \emph{theory-grounded types};
they are not user clusters or fitted population proportions.  The public data
deposit contains tweet identifiers \parencite{debnath2023dataset}.  The thesis
did not hydrate and recluster the 814,924 tweets and therefore does not claim to
replicate Debnath et al.'s analysis, infer their social graph, or construct a
Twitter digital twin.

This separation supports a defensible use of the case.  The observational study
supplies the empirical phenomenon---chemtrails-related spillover, affect, and
toxicity---while the simulation supplies counterfactual control over seeds,
identity composition, and network structure.  Agreement between the two would
be evidence of external plausibility, not proof that the simulated agents are
the users in Debnath et al.'s data.

\section{Network topology, homophily, and echo}
\label{sec:lit-topology}

Topology changes who can expose whom and how many independent or reinforcing
paths connect groups.  Erd\H{o}s--R\'enyi graphs provide a random-graph reference
\parencite{erdos1960evolution}; Watts--Strogatz networks combine local clustering
with short paths \parencite{watts1998collective}; and preferential attachment
generates heavy-tailed degree distributions and hubs
\parencite{barabasi1999emergence}.  These models are not interchangeable proxies
for real social media.  They are controlled structural hypotheses whose degree
distribution, clustering, path length, and seed position must be reported.

The content being transmitted also changes which topology should facilitate
diffusion.  \textcite{centola2007complex} distinguish simple contagion, for which
one contact may suffice, from complex contagion requiring reinforcement from
multiple neighbours.  Long weak ties can accelerate the former but fail to
provide sufficiently wide bridges for the latter.  Firestorm participation and
conspiracy reinterpretation may involve social affirmation and are therefore
not safely assumed to behave like infection over a single edge.

Agent composition and topology are joined by homophily.  Similarity structures
social ties and thereby the information, attitudes, and interactions available
to people \parencite{mcpherson2001homophily}.  Empirically, selective exposure and
homogeneous interaction can support polarised information communities:
\textcite{delvicario2016echo} relate emotional contagion and group polarisation
in Facebook communities, while \textcite{tornberg2018echo} models viral
misinformation as a complex contagion.  A high-level ``echo chamber'' label,
however, is not a measurement.  This thesis therefore distinguishes treatment
variables (network generator and persona assignment) from observables such as
edge homophily and community separation.  Modularity compares within-group
edge density with a degree-preserving null expectation
\parencite{newman2006modularity}; it indicates structural segregation, not the
truth or extremity of group content.

\section{Homogeneous and heterogeneous agent populations}
\label{sec:lit-heterogeneity}

Homogeneous-agent models isolate a response rule and simplify attribution, but
they can turn an imposed persona into a population-wide mechanism.
Heterogeneous populations permit cross-cutting exposure, specialist roles, and
different propensities to reinterpret or resist a message.  At the same time,
heterogeneity is not intrinsically realistic.  Results depend on which types are
included, their prevalence, their network positions, and whether prompts create
genuine behavioural variation or merely labels.

The immediate methodological predecessor is the auditor--node system of
\textcite{maurya2025simulating}, presented at the LASS workshop at ACM CIKM
2025.  It compares persona-homogeneous and persona-heterogeneous linear rewrite
chains and introduces the discrete Misinformation Index (MI) and branch-level
Misinformation Propagation Rate (MPR).  Those general-news linear-chain findings
are \emph{prior work}, not results of this thesis.  The present work retains the
auditor--node idea but changes the empirical object to geoengineering discourse,
uses Debnath-grounded profiles, and introduces non-chain network treatments.

The broader literature demonstrates both the promise and the unresolved status
of LLM societies.  Social Simulacra populates prototypes of online communities
\parencite{park2022simulacra}, while Generative Agents combines memory,
reflection, and planning to create believable local interactions
\parencite{park2023generative}.  S$^3$ simulates information, attitude, and
emotion propagation against social-network data \parencite{gao2023s3}.  OASIS
adds platform actions, recommendation mechanisms, dynamic networks, and very
large populations \parencite{yang2024oasis}; AgentSociety similarly prioritises
large-scale social environments \parencite{piao2025agentsociety}.  These systems
expand scale and behavioural surface area, but scale does not by itself validate
an agent as a model of a person.

Most directly, \textcite{li2026diffusion} condition LLM-agent forwarding on
psychological traits and compare diffusion under different network structures,
including an empirical Higgs community.  Their topology-sensitive forwarding
study is close to this thesis, but the focal state remains sharing and diffusion
rather than audited claim-by-claim mutation in geoengineering text.  The
literatures are complementary: one asks whether an item is forwarded; this
thesis asks how the item's factual content changes as a function of identity
mix and topology.

Persona conditioning also creates a validity problem.  Language models can
reproduce some subgroup response patterns when conditioned on demographic
backstories \parencite{argyle2023one}, and they can recover qualitative patterns
from selected economic experiments \parencite{horton2023silicus}.  Such results
support LLMs as inexpensive theory probes, not synthetic replacements for human
participants.  Training-data contamination, prompt sensitivity, stochasticity,
stereotype amplification, limited within-group variance, and obedience to
experimenter instructions all threaten behavioural interpretation.  In this
thesis, ``homogeneous'' and ``heterogeneous'' consequently describe prompt
conditions in a computational experiment.  They do not denote representative
human populations.

\section{Measurement validity: discrete and continuous instruments}
\label{sec:lit-validity}

The thesis's central dependent variable is produced by an instrument, not
observed directly.  Construct validity concerns whether evidence supports the
interpretation and use of a score.  It cannot be established by naming a
quantity ``misinformation.''  \textcite{cronbach1955construct} require a measure
to sit in a nomological network connecting a construct to observable
consequences, and \textcite{messick1995validity} treats validity as an integrated
argument about score meaning, generalisation, external relations, and
consequences.  For MI, relevant evidence would include item coverage, human
agreement, expected relations with known distortions, stability across judges,
and sensitivity to irrelevant wording.

QA-based factuality evaluation supplies a methodological warrant.
\textcite{fabbri2022qafacteval} show that question answering can assess factual
consistency in summarisation more directly than lexical overlap.  LLM judges can
also agree substantially with human preferences, but exhibit position,
verbosity, and self-enhancement biases \parencite{zheng2023judge}.  Neither result
validates this thesis's exact instrument: fixed binary questions about
geoengineering, answered by an LLM judge, are a different construct and use
case.

The distinction between discrete and continuous scoring must therefore remain
visible.  In discrete MI, each item is classified as recovered, missing, or
contradicted and the item outcomes are aggregated into a bounded score.
This format is inspectable and corresponds to explicit propositions, but it
collapses uncertainty and treats all items as equally diagnostic.  A continuous
instrument can retain graded evidence or confidence and may detect smaller
changes, but its numeric precision is not automatically substantive precision.
It requires calibration and a defensible interpretation of distances between
scores.

The statistical literature cautions that converting a genuinely continuous
quantity into categories discards information, reduces power, and creates
artificial discontinuities around a cut point
\parencite{maccallum2002dichotomization,altman2006cost}.  That warning applies to
the thesis's severity bands: the boundary at MI${}>3$ is an operational decision
inherited from prior work, not a psychometrically established natural kind.
Conversely, the warning does not prove that a continuous LLM score is valid.
Discrete and continuous MPR are different instruments; they should not be pooled
or treated as interchangeable observations merely because both have numeric
output.

Shared generator and judge families add common-method dependence.  A fluent
rewrite may be rewarded because it resembles the judge's preferred formulation,
and a failed parse can masquerade as low misinformation if the implementation
fails open.  Credible validation therefore requires frozen test cases,
fail-closed parsing, deterministic judging where possible, manual inspection,
human coding of a sample, and agreement reported at the decision boundary used
for the substantive claim.  Until such evidence exists, MI and MPR are best
described as \emph{simulation audit scores}, not validated measures of human
misinformation belief.

\section{Validation of generative social simulation}
\label{sec:lit-simulation-validation}

Verification asks whether the program implements the intended model; validation
asks whether the model is adequate for its intended purpose.
\textcite{sargent2013verification} separates data validity, conceptual-model
validity, computerised-model verification, and operational validity.  For
agent-based social simulation, \textcite{richiardi2006protocol} further stress
transparent model description, sensitivity analysis, replication, and a
validation test matched to the claim.  \textcite{windrum2007validation} show that
empirical validation can involve several strategies and that no single fitted
trajectory settles model adequacy.

Applied here, validation has at least four levels.  First, software checks must
establish that topology, persona assignment, intervention timing, and auditor
aggregation behave as specified.  Second, internal validation asks whether
changing a theoretically relevant input produces stable, interpretable changes
rather than parser artefacts or dead cascades.  Third, output validation compares
patterns with held-out empirical regularities, such as plausible cascade
structure or Debnath's discourse themes.  Fourth, measurement validation asks
whether humans and independent instruments support the MI interpretation.
Face-plausible agent prose is evidence for none of these levels by itself.

The intended use defines the attainable claim.  A small, replicated, audited
model may function as a mechanism probe: it can test whether a proposed
combination of topology and identity is sufficient to produce a pattern under
specified prompts.  Without calibrated agent distributions and reconstructed
empirical cascades, it cannot estimate the prevalence or causal effect of that
mechanism among Twitter users.  Failed or near-empty cascades are execution or
design outcomes to diagnose, not evidence that a topology prevents
misinformation.  Stochastic LLM output also requires multiple seeds or models
before a directional contrast can be considered stable.

\section{Synthesis and research gap}
\label{sec:lit-gap}

The literature leaves a precise intersection underdeveloped.  Firestorm theory
identifies structural conditions for rapid, affective outrage but does not
provide a generative, claim-level model \parencite{pfeffer2014firestorms}.
Network diffusion research explains activation, reinforcement, and cascade
shape but usually treats content as a binary or fixed state
\parencite{kempe2003influence,centola2007complex}.  Debnath et al. provide a
large observational map of geoengineering conspiracy spillovers but not
controlled counterfactual societies \parencite{debnath2023conspiracy}.  LLM
societies generate heterogeneous behaviour at increasing scale, yet their
misinformation evaluations commonly emphasise sharing, aggregate plausibility,
or believability rather than a validated trajectory of proposition-level change.
Finally, the prior auditor--node study supplies MI/MPR and homogeneous versus
heterogeneous chains, but not a geoengineering-grounded network experiment
\parencite{maurya2025simulating}.

The gap filled by this thesis is therefore \emph{not} a new general-purpose agent
architecture and not a replication of an 814,924-tweet study.  It is a
theory-mapped, auditable test of how network topology and Debnath-grounded
identity composition jointly condition semantic misinformation trajectories in
geoengineering discourse.  It connects (i) an explicit seven-to-six Pfeffer
operationalisation, (ii) controlled homogeneous and heterogeneous LLM-agent
societies, and (iii) discrete and continuous claim-recovery instruments whose
validity limits remain explicit.  The empirical contribution is the resulting
climate-domain comparison and its validation evidence; the methodological
contribution is a separation of transmission, semantic mutation, firestorm
timing, and instrument choice.  This design addresses a question not answered
by any one parent literature: under fixed scientific seeds, are apparent
firestorm thresholds robust to who occupies the network, how those agents are
connected, and how factual degradation is measured?
